\section{Homework 4}

\begin{exercise}
    [P34, Ex2] \(M=\mathfrak{sl}(3,F)\) contains a copy of \(L\) in its upper left-hand \(2\times 2\) position. Write \(M\) as a direct sum of irreducible \(L\)-submodules, where \(M\) is viewed as an \(L\)-module via the adjoint representation:
    \[
        V(0)\oplus V(1)\oplus V(1)\oplus V(2).
    \]
\end{exercise}

\begin{proof}
    Let
    \[
    L=\mathfrak{sl}_2(F)
     =
    F e\oplus Fh\oplus Ff
    \]
    be embedded in the upper left-hand \(2\times 2\) corner of
    \[
    M=\mathfrak{sl}_3(F),
    \]
    where
    \[
    e=E_{12},\qquad f=E_{21},\qquad h=E_{11}-E_{22}.
    \]
    We view \(M\) as an \(L\)-module by the adjoint action:
    \[
    x\cdot m=[x,m]\qquad (x\in L,\;m\in M).
    \]
    
    Set
    \[
    z=E_{11}+E_{22}-2E_{33}.
    \]
    Then \(z\in \mathfrak{sl}_3(F)\), and \(z\) commutes with \(e,f,h\), because
    it is scalar on the upper left-hand \(2\times 2\) block. Hence
    \[
    M_0:=Fz
    \]
    is a trivial \(L\)-submodule, so
    \[
    M_0\cong V(0).
    \]
    
    Next, the embedded copy of \(L\),
    \[
    M_2:=F E_{12}\oplus F(E_{11}-E_{22})\oplus F E_{21},
    \]
    is stable under the adjoint action of \(L\). It is just the adjoint module of
    \(\mathfrak{sl}_2(F)\). Since \(E_{12}=e\) is a highest-weight vector of weight
    \(2\), this submodule is irreducible and
    \[
    M_2\cong V(2).
    \]
    
    Now consider
    \[
    M_1^+:=F E_{13}\oplus F E_{23}.
    \]
    Using
    \[
    [E_{ij},E_{kl}]=\delta_{jk}E_{il}-\delta_{li}E_{kj},
    \]
    we get
    \[
    [h,E_{13}]=E_{13},\qquad [h,E_{23}]=-E_{23},
    \]
    \[
    [e,E_{13}]=0,\qquad [e,E_{23}]=E_{13},
    \]
    and
    \[
    [f,E_{13}]=E_{23},\qquad [f,E_{23}]=0.
    \]
    Thus \(E_{13}\) is a highest-weight vector of weight \(1\), and
    \[
    M_1^+\cong V(1).
    \]
    
    Similarly define
    \[
    M_1^-:=F E_{32}\oplus F E_{31}.
    \]
    We have
    \[
    [h,E_{32}]=E_{32},\qquad [h,E_{31}]=-E_{31},
    \]
    \[
    [e,E_{32}]=0,\qquad [e,E_{31}]=-E_{32},
    \]
    and
    \[
    [f,E_{32}]=-E_{31},\qquad [f,E_{31}]=0.
    \]
    Hence \(E_{32}\) is again a highest-weight vector of weight \(1\), and,
    after changing the sign of one basis vector if desired, this submodule is
    isomorphic to \(V(1)\):
    \[
    M_1^-\cong V(1).
    \]
    
    Finally,
    \[
    \dim M_0+\dim M_1^++\dim M_1^-+\dim M_2
    =
    1+2+2+3=8=\dim \mathfrak{sl}_3(F),
    \]
    and the displayed subspaces are spanned by disjoint subsets of the standard
    matrix-unit basis together with the diagonal trace-zero elements. Therefore their
    sum is direct and equals \(M\). Hence
    \[
    M
    =
    M_0\oplus M_1^+\oplus M_1^-\oplus M_2
    \cong
    V(0)\oplus V(1)\oplus V(1)\oplus V(2).
    \]
    \end{proof}

\begin{exercise}
    [P34, Ex6] Decompose the tensor product of the two \(L\)-modules \(V(3)\) and \(V(7)\) into the sum of irreducible submodules:
    \[
        V(4)\oplus V(6)\oplus V(8)\oplus V(10).
    \]
    Try to develop a general formula for the decomposition of \(V(m)\otimes V(n)\).
\end{exercise}

\begin{proof}
    For \(r\geq 0\), the irreducible \(\mathfrak{sl}_2(F)\)-module \(V(r)\) has
    weights
    \[
    r,\; r-2,\; r-4,\;\ldots,\;-r,
    \]
    each with multiplicity \(1\). Hence its formal character is
    \[
    \operatorname{ch}V(r)=q^r+q^{r-2}+\cdots+q^{-r}.
    \]
    
    Assume first that \(0\leq m\leq n\). In \(V(m)\otimes V(n)\), the weights are
    sums of weights from \(V(m)\) and \(V(n)\). Thus the multiplicity of the weight
    \[
    m+n-2s
    \]
    is
    \[
    c_s
    =
    \#\{(i,j)\mid 0\leq i\leq m,\;0\leq j\leq n,\; i+j=s\}.
    \]
    Therefore
    \[
    c_s=
    \begin{cases}
    s+1, & 0\leq s\leq m,\\
    m+1, & m\leq s\leq n,\\
    m+n-s+1, & n\leq s\leq m+n.
    \end{cases}
    \]
    
    On the other hand, consider
    \[
    \bigoplus_{k=0}^{m} V(m+n-2k).
    \]
    The weight \(m+n-2s\) occurs in \(V(m+n-2k)\) precisely when
    \[
    k\leq s\leq m+n-k.
    \]
    Thus its multiplicity in the above direct sum is
    \[
    \#\{k\mid 0\leq k\leq m,\; k\leq s,\; k\leq m+n-s\}
    =
    \min\{m,s,m+n-s\}+1,
    \]
    which gives exactly the same three cases as \(c_s\). Hence
    \[
    \operatorname{ch}\bigl(V(m)\otimes V(n)\bigr)
    =
    \operatorname{ch}\left(\bigoplus_{k=0}^{m}V(m+n-2k)\right).
    \]
    Since finite-dimensional \(\mathfrak{sl}_2(F)\)-modules are completely reducible
    and irreducibles are determined by their highest weights, we obtain
    \[
    V(m)\otimes V(n)
    \cong
    \bigoplus_{k=0}^{m}V(m+n-2k),
    \qquad 0\leq m\leq n.
    \]
    
    Equivalently, for arbitrary \(m,n\geq 0\),
    \[
    V(m)\otimes V(n)
    \cong
    \bigoplus_{k=0}^{\min\{m,n\}} V(m+n-2k).
    \]
    
    Applying this with \(m=3\) and \(n=7\), we get
    \[
    V(3)\otimes V(7)
    \cong
    V(10)\oplus V(8)\oplus V(6)\oplus V(4).
    \]
    Reordering the direct summands,
    \[
    V(3)\otimes V(7)
    \cong
    V(4)\oplus V(6)\oplus V(8)\oplus V(10).
    \]
    \end{proof}

\begin{exercise}
    [P41, Ex1] If \(L\) is a classical linear Lie algebra of type \(A_\ell\), \(B_\ell\), \(C_\ell\), or \(D_\ell\) (see (1.2)), prove that the set of all diagonal matrices in \(L\) is a maximal toral subalgebra, of dimension \(\ell\). (Cf. Exercise 2.8.)
\end{exercise}

\begin{proof}
    We use the standard matrix realizations of the classical Lie algebras.
    
    For type \(A_\ell\),
    \[
    L=\mathfrak{sl}_{\ell+1}(F),
    \]
    and the diagonal matrices in \(L\) are
    \[
    H
    =
    \left\{
    \operatorname{diag}(a_1,\ldots,a_{\ell+1})
    \;\middle|\;
    \sum_{i=1}^{\ell+1}a_i=0
    \right\}.
    \]
    Hence
    \[
    \dim H=\ell.
    \]
    
    For type \(B_\ell\), in the standard realization of \(\mathfrak{o}_{2\ell+1}(F)\),
    the diagonal matrices in \(L\) have the form
    \[
    H
    =
    \left\{
    \operatorname{diag}(0,a_1,\ldots,a_\ell,-a_1,\ldots,-a_\ell)
    \right\}.
    \]
    Thus
    \[
    \dim H=\ell.
    \]
    
    For types \(C_\ell\) and \(D_\ell\), in the standard realizations
    \[
    C_\ell:\quad L=\mathfrak{sp}_{2\ell}(F),
    \qquad
    D_\ell:\quad L=\mathfrak{o}_{2\ell}(F),
    \]
    the diagonal matrices in \(L\) have the form
    \[
    H
    =
    \left\{
    \operatorname{diag}(a_1,\ldots,a_\ell,-a_1,\ldots,-a_\ell)
    \right\}.
    \]
    Again,
    \[
    \dim H=\ell.
    \]
    
    In each case, \(H\) is abelian because diagonal matrices commute. Moreover, for
    a diagonal matrix \(h=\operatorname{diag}(d_1,\ldots,d_N)\), one has
    \[
    [h,E_{ij}]=(d_i-d_j)E_{ij}.
    \]
    Thus \(\operatorname{ad}h\) is diagonalizable on the ambient matrix algebra, and
    therefore also on the invariant subspace \(L\). Hence every element of \(H\) is
    semisimple under the adjoint representation. Therefore \(H\) is a toral subalgebra.
    
    It remains to prove maximality. We show that
    \[
    C_L(H)=H.
    \]
    Let \(V\) be the natural module. In type \(A_\ell\), the weights of \(H\) on \(V\)
    are
    \[
    \varepsilon_1,\ldots,\varepsilon_{\ell+1},
    \]
    restricted to \(\sum a_i=0\), and they are distinct. In type \(B_\ell\), the weights
    are
    \[
    0,\;\pm\varepsilon_1,\ldots,\pm\varepsilon_\ell,
    \]
    and in types \(C_\ell,D_\ell\), they are
    \[
    \pm\varepsilon_1,\ldots,\pm\varepsilon_\ell.
    \]
    These weights are distinct. Hence any element of \(\mathfrak{gl}(V)\) commuting
    with every element of \(H\) preserves each one-dimensional weight space of \(V\).
    Therefore it is diagonal in the standard basis. Consequently,
    \[
    C_{\mathfrak{gl}(V)}(H)
    \]
    is the algebra of diagonal matrices. Intersecting with \(L\), we get exactly the
    diagonal matrices lying in \(L\), namely
    \[
    C_L(H)=H.
    \]
    
    Now let \(T\) be a toral subalgebra of \(L\) containing \(H\). Since toral subalgebras
    are abelian, every element of \(T\) commutes with every element of \(H\). Hence
    \[
    T\subseteq C_L(H)=H.
    \]
    Thus \(T=H\), proving that \(H\) is maximal toral.
    
    Therefore, in each of the classical types
    \[
    A_\ell,\quad B_\ell,\quad C_\ell,\quad D_\ell,
    \]
    the set of all diagonal matrices in \(L\) is a maximal toral subalgebra of dimension
    \(\ell\).
    \end{proof}

\begin{exercise}
    [P41, Ex10] Prove that no four, five, or seven dimensional semisimple Lie algebras exist.
\end{exercise}

\begin{proof}
    Let \(L\) be a nonzero semisimple Lie algebra over \(F\). Choose a maximal toral
    subalgebra \(H\), and let
    \[
    \ell=\dim H
    \]
    be the rank of \(L\). Let \(\Phi\subset H^*\) be the corresponding root system.
    Then the root-space decomposition gives
    \[
    L
    =
    H\oplus \bigoplus_{\alpha\in\Phi} L_\alpha,
    \]
    where
    \[
    \dim L_\alpha=1
    \]
    for every \(\alpha\in\Phi\). Therefore
    \[
    \dim L=\ell+|\Phi|.
    \]
    
    Since \(\Phi\) is a root system, \(\alpha\in\Phi\) implies \(-\alpha\in\Phi\), and
    \(0\notin\Phi\). Thus \(|\Phi|\) is even, so
    \[
    \dim L\equiv \ell \pmod 2.
    \]
    Also, since \(\Phi\) spans \(H^*\), we may choose \(\ell\) linearly independent
    roots
    \[
    \alpha_1,\ldots,\alpha_\ell.
    \]
    Then
    \[
    \pm\alpha_1,\ldots,\pm\alpha_\ell
    \]
    are \(2\ell\) distinct roots, so
    \[
    |\Phi|\geq 2\ell.
    \]
    Hence
    \[
    \dim L=\ell+|\Phi|\geq 3\ell.
    \]
    
    If \(\ell=1\), then \(\Phi\) is a rank-one root system. Since no scalar multiples of
    a root occur except its negative, we must have
    \[
    \Phi=\{\alpha,-\alpha\}.
    \]
    Therefore
    \[
    \dim L=\ell+|\Phi|=1+2=3.
    \]
    
    We now rule out the dimensions \(4,5,7\).
    
    If \(\dim L=4\), then
    \[
    4\equiv \ell \pmod 2,
    \]
    so \(\ell\) is even. Since \(L\neq 0\), we have \(\ell>0\), hence \(\ell\geq 2\).
    But then
    \[
    \dim L\geq 3\ell\geq 6,
    \]
    contradicting \(\dim L=4\).
    
    If \(\dim L=5\), then \(\ell\) is odd. Since
    \[
    5=\dim L\geq 3\ell,
    \]
    we must have \(\ell=1\). But rank one gives
    \[
    \dim L=3,
    \]
    contradicting \(\dim L=5\).
    
    If \(\dim L=7\), then \(\ell\) is odd. Since
    \[
    7=\dim L\geq 3\ell,
    \]
    we again must have \(\ell=1\). But then
    \[
    \dim L=3,
    \]
    contradicting \(\dim L=7\).
    
    Thus no semisimple Lie algebra of dimension \(4\), \(5\), or \(7\) exists.
    \end{proof}
